\chapter*{ABSTRACT}
\onehalfspacing
Gene expression microarray data typically comprise measurements for tens of thousands of probes recorded on a relatively small number of biological samples. Such high-dimensional datasets present significant challenges for machine learning models, including an increased risk of overfitting, higher computational cost, and reduced generalization performance. Feature selection techniques are commonly used to mitigate these challenges by identifying a smaller, more informative subset of probes for classification. This study investigated the effect of filter, wrapper, and embedded feature selection methods on the classification performance of a linear Support Vector Machine (SVM) applied to two Affymetrix microarray gene-expression datasets obtained from the Gene Expression Omnibus (GEO): GSE19804 (120 samples: 60 lung-cancer, 60 normal; class-balanced) and GSE42568 (121 samples: 104 breast-cancer, 17 normal; severely imbalanced). Both datasets contained 54,675 probes; no probe-to-gene annotation was performed, and reported feature counts therefore refer to probes rather than uniquely mapped genes. Seven feature-selection variants; Welch's $t$-test, the ANOVA $F$-test, Benjamini-Hochberg FDR-ranked filtering, SVM Recursive Feature Elimination, Random Forest importance, XGBoost gain importance, and LASSO were compared against a full-feature baseline using 5-fold stratified cross-validation, with all scaling and feature-selection steps fitted strictly within each training fold to avoid data leakage.

On the balanced GSE19804 dataset, several feature-selection methods matched the full-feature baseline (95.00\% accuracy, MCC $=0.9013$), while the XGBoost-importance selector achieved the highest mean performance (97.50\% accuracy, MCC $=0.9530$); SVM-RFE was the only method that reduced performance relative to the baseline (92.50\% accuracy, MCC $=0.8554$). On the severely imbalanced GSE42568 dataset, none of the evaluated methods produced a reliable classifier: most configurations collapsed to majority-class prediction (0\% specificity), and even the best-performing configuration, BH-FDR-ranked selection retaining two probes, achieved only MCC $=0.0269$, close to chance. These findings indicate that the effect of feature selection on classification performance is method- and dataset-dependent rather than uniformly beneficial, and that severe class imbalance combined with a very small minority class can prevent reliable classification even after aggressive dimensionality reduction.
On the balanced GSE19804 dataset, several feature-selection methods matched the full-feature baseline (95.00\% accuracy, MCC $=0.9013$), while the XGBoost-importance selector achieved the highest mean performance (97.50\% accuracy, MCC $=0.9530$); SVM-RFE was the only method that reduced performance relative to the baseline (92.50\% accuracy, MCC $=0.8554$). On the severely imbalanced GSE42568 dataset, none of the evaluated methods produced a reliable classifier: most configurations collapsed to majority-class prediction (0\% specificity), and even the best-performing configuration---BH-FDR-ranked selection retaining two probes---achieved only MCC $=0.0269$, close to chance. These findings indicate that the effect of feature selection on classification performance is method- and dataset-dependent rather than uniformly beneficial, and that severe class imbalance combined with a very small minority class can prevent reliable classification even after aggressive dimensionality reduction.\footnote{\textbf{[Action item before submission:]} the figures above were computed before a fold-locality fix to the standardization step (see Chapter~3); they should be refreshed from the corrected pipeline before final submission, though only a small numerical shift is expected.}
